\graphicspath{{conclusion/fig/}}

\mychapter{Summary and conclusions}{Summary and Conclusions}{}\label{chap:conclusion}

\section{Summary}
\label{sec:conc_summary}

This project investigated whether a CNN can determine the state of the aerobic zones and secondary clarifiers of a WWTW from Google Earth images. Five models were trained for each component and evaluated on a chronological split, which represents a known facility, and on a leave-one-facility-out split, which represents a new facility (Chapter~\ref{chap:experimental_design}).

\section{Conclusions}
\label{sec:conc_conclusions}

At a known facility, the clarifier models reach a mean AUC of 0.921 to 0.949 and the aerobic zone models 0.860 to 0.938 (Table~\ref{tab:res_multiclass}), so the state of both components can be determined from satellite images. At a new facility, the clarifier models perform about as well, but the macro recall of the aerobic zone models decreases by 12.2 to 30.5 percentage points. An aerobic zone model should therefore only be used at a facility whose images are part of the training data, which is the situation in the intended practical application of the system. As binary classifiers, the models reach or come close to the target Not-OK recall of 95\% at a known facility, but not at a new facility (Table~\ref{tab:res_binary}), so a threshold chosen on other facilities cannot be relied on. The Grad-CAM heatmaps mostly lie on the features that the labels are based on, but for some clarifier classes they lie on the rim, which carries no information about the state of the clarifier.

\section{Limitations}
\label{sec:conc_limitations}

The aerobic zone labels were not reviewed, so some of the aerobic zone mistakes may be caused by the labels. The aerobic zones are also labelled at only three facilities, so their leave-one-facility-out results are based on three test sets. Finally, the hyperparameters were tuned only with Northern Works held out for the clarifiers and Atlantis held out for the aerobic zones, and the same values were used in every experiment, so they may not be optimal for the other facilities.

\section{Future Work}
\label{sec:conc_future}

The next step is to implement the dashboard, where new images are classified, components classified as Not-OK are relabelled by an operator, and the relabelled images are added to the training data. The classification may also be improved with satellite data that has more spectral bands, such as near-infrared, or a higher resolution than the RGB images from Google Earth.
